%! Author = zhiweizhang
%! Date = 10.06.25

% Preamble
\chapter{Thesis Structure}\label{chapter:thesis-structure}

This thesis is organized as follows:

\begin{itemize}
  \item \textbf{Chapter 2: Literature Review} – We examine existing research and methodologies related to Nazi
  symbol detection and classification, highlighting key findings and identifying gaps that our work aims to address.

  \item \textbf{Chapter 3: Methodology} – This chapter details our approach, including data collection, dataset
  composition, preprocessing techniques, and the machine learning algorithms used for classification. We also
  discuss model training, evaluation metrics, and the rationale behind our chosen methods.

  \item \textbf{Chapter 4: Ethical Considerations} – We reflect on the ethical dimensions of collecting and analyzing
  sensitive data, including privacy concerns, potential misuse of models, and the responsibilities involved in automated
  hate symbol detection.

  \item \textbf{Chapter 5: Results} – We present the outcomes of our experiments, providing quantitative and
  qualitative analyses of the classification system’s performance.

  \item \textbf{Chapter 6: Discussion and Analysis} – We interpret the results, discuss their implications, and
  compare our findings with previous research. We also explore the system's limitations and potential improvements.

  \item \textbf{Chapter 7: Model Deployment and API Service} – This chapter describes the design and implementation
  of an API service for deploying the best-performing models, including the system architecture, CI/CD pipeline, and
  available endpoints.

  \item \textbf{Chapter 8: Conclusion and Future Work} – This final chapter summarizes our contributions, outlines
  the key takeaways, and suggests directions for future research in Nazi symbol classification.
\end{itemize}